\chapter{Sidon sets and elementary foundations}

\begin{definition}[Sidon set]
\lean{IsSidonSet}\leanok\label{def:sidon_set}
For an additive commutative monoid $M$, a set $S\subseteq M$ is Sidon
if $a+b=c+d$ with $a,b,c,d\in S$ implies $\{a,b\}=\{c,d\}$.
\end{definition}

\begin{theorem}[Empty set]
\lean{empty_is_sidon}\leanok\uses{def:sidon_set}
The empty set is Sidon.
\end{theorem}
\begin{proof}
Any four membership hypotheses include an element of the empty set,
a contradiction. The defining implication follows.
\end{proof}

\begin{theorem}[Singleton]
\lean{singleton_is_sidon}\leanok\uses{def:sidon_set}
Every singleton is Sidon.
\end{theorem}
\begin{proof}
Membership forces each element to equal the singleton's element.
Substitution makes the required unordered-pair equality reflexive.
\end{proof}

\begin{theorem}[Hereditary property]
\lean{sidon_subset_is_sidon}\leanok\uses{def:sidon_set}
If $S$ is Sidon and $T\subseteq S$, then $T$ is Sidon.
\end{theorem}
\begin{proof}
Use the inclusion on the four membership hypotheses, then apply
the defining implication for $S$ to the same sum equality.
\end{proof}

\begin{lemma}[Unique nonzero differences]
\lean{IsSidonSet.eq_of_sub_eq_sub,IsSidonSet.iff_unique_nonzero_difference}
\leanok\label{found:diff}\uses{def:sidon_set}
In an additive commutative group, $S$ is Sidon if and only if
$a,b,c,d\in S$, $a\ne b$, $a-b=c-d$ imply $a=c$ and $b=d$.
\end{lemma}
\begin{proof}
Rewrite the difference equality as $a+d=c+b$. The Sidon condition
gives $\{a,d\}=\{c,b\}$, so $a=c$ or $a=b$. Exclude the second
case and cancel in the difference equality to obtain $b=d$.
Conversely, suppose $a+b=c+d$. If $a=c$, cancellation gives $b=d$.
Otherwise rewrite as $a-c=d-b$ and use difference uniqueness on
$(a,c)$ and $(d,b)$ to obtain $a=d$ and $c=b$.
Both cases give the required unordered-pair equality.
Finally $a-b\ne0$ is equivalent to $a\ne b$, identifying this with
the paper's unique-nonzero-difference definition. No torsion assumption
is used.
\end{proof}

\begin{theorem}[Finite group bound]
\lean{sidon_set_size_bound}\leanok\label{found:finite}
\uses{def:sidon_set,found:diff}
For Sidon $S$ in a finite additive commutative group $M$,
$|S|(|S|-1)\leq|M|$. In fact $|S|(|S|-1)\leq|M|-1$.
The first declaration gives the weaker bound; the refinement is checked as
\texttt{IsSidonSet.card\_sq\_sub\_card\_le\_card\_sub\_one}.
\end{theorem}
\begin{proof}
$S^2$ has $m^2$ pairs and its diagonal has $m$ pairs, so its
off-diagonal has $m^2-m=m(m-1)$ pairs, also when $m=0$.
The difference map on that off-diagonal is injective by
Lemma~\ref{found:diff}, and takes values in $M\setminus\{0\}$.
The latter has $|M|-1$ elements, since $M$ contains zero.
Finite injectivity proves the refined bound and hence the weaker one.
\end{proof}

\begin{corollary}[Cyclic groups]
\lean{sidon_set_size_bound_zmodn}\leanok\uses{found:finite}
For $n\geq1$ and Sidon $S\subseteq\mathbb Z/n\mathbb Z$,
$|S|(|S|-1)\leq n$.
\end{corollary}
\begin{proof}
For $n\ne0$, $\mathbb Z/n\mathbb Z$ has $n$ elements.
Substitute this cardinality into the finite group bound.
\end{proof}

\begin{lemma}[Affine images over the integers]
\lean{IsSidonSet.image_affine,Sidon.card_image_affine}\leanok
\label{found:affine}\uses{def:sidon_set}
For integers $u\ne0$ and $v$, the set $uS+v$ is Sidon whenever
$S\subseteq\mathbb Z$ is Sidon. Finite such images have cardinality $|S|$.
\end{lemma}
\begin{proof}
Choose four preimages for a sum equality in the image. Expanding and
canceling $2v$ and then the nonzero factor $u$ gives a sum equality in
$S$. Apply the Sidon condition and take the image of the resulting
unordered-pair equality. If $ua+v=ub+v$, cancel $v$ and $u$ to get
$a=b$. Thus the map is injective and preserves finite cardinality.
\end{proof}

\begin{lemma}[Crude interval bound]\label{found:crude}\uses{found:diff}
\lean{IsSidonSet.card_le_two_mul_sqrt}\leanok
A Sidon subset of an integer interval of length $L\geq1$ has
at most $2\sqrt L$ elements. At $L=0$ it is empty.
\end{lemma}
\begin{proof}
The $m(m-1)$ distinct nonzero ordered differences lie among
$\{-(L-1),\ldots,-1,1,\ldots,L-1\}$, giving
$m(m-1)\leq2(L-1)$.
If $m>2\sqrt L$, then $\sqrt L\geq1$ gives
$m(m-1)>2\sqrt L(2\sqrt L-1)\geq2L>2(L-1)$,
a contradiction. An interval of length zero is empty.
This bound alone suffices for the Fourier theorem's boundary terms.
\end{proof}

\begin{proposition}[Odd construction, paper Proposition 1.13]
\label{found:odd}\uses{found:affine}
Let $L=\lfloor(N-1)/2\rfloor$ and $S_0\subseteq[L]$ be Sidon.
Then $S=2S_0+1\subseteq[N]$ is Sidon, $|S|=|S_0|$, and it has
no solution to $x_1+x_2+x_3+x_4=x_5$.
If $|S_0|=\sqrt L(1+o(1))$, then
$|S|=\sqrt{N/2}(1+o(1))$ as $N\to\infty$.
\end{proposition}
\begin{proof}
Use the affine lemma. For $1\leq x\leq L$,
$1\leq2x+1\leq2L+1\leq N$. Every image element is odd, but a
sum of four odd integers is even. Reduction modulo two rules out the
equation. The floor bound gives $L/N\to1/2$; continuity of square
root gives $\sqrt{L/N}\to1/\sqrt2$. Multiply by
$|S_0|/\sqrt L\to1$. The existence of asymptotically optimal $S_0$
is the same external construction input as in
Corollary~\ref{target:largest}.
\end{proof}
